% ===========================================================================
%  CAPITULO 7 — Consolidação para o Termo de Referência
% ===========================================================================
\chapter{Consolidação para o Termo de Referência}
\label{cap:consolidacao}

\section{O que o TR pode afirmar com segurança}

Itens com evidência no código e em testes, que podem ser exigidos como
requisito sem risco de a solução atual não atender:

\begin{itemize}
  \item Acesso 100\% web, sem instalação, com interface responsiva.
  \item Carga cadastral no padrão OneRoster 1.2 (CSV), por interface e API,
    com importação parcial, relatório de erros e carga idempotente.
  \item Seis perfis com autorização aplicada no servidor e isolamento lógico
    por escola; isolamento físico por contrato.
  \item Ficha e guia de aula por missão, com objetivos, habilidades BNCC,
    roteiro e rubrica.
  \item Personalização por IA com aprovação obrigatória do professor antes do
    uso, sem alterar a estrutura pedagógica nem os critérios de avaliação.
  \item Avaliação por habilidade com regras no servidor, histórico de rodadas
    e revisão pelo professor, que prevalece sobre a classificação
    automática.
  \item Proteção do texto do estudante: nada de identificação enviada à IA;
    textos livres fora de painéis e registros técnicos.
  \item Progresso salvo no servidor a cada interação, com retomada.
\end{itemize}

\section{O que o TR não deve afirmar ainda}

\begin{limitacao}
\begin{itemize}
  \item Ganhos de aprendizagem, engajamento ou aceitação (P10 sem relatório).
  \item Capacidade comprovada de carga (nenhum teste executado, T09).
  \item Conformidade declarada com WCAG AA (falta a avaliação manual, P09).
  \item Login federado por OIDC ou SAML em operação (T02).
  \item Duplicação ou compartilhamento de missões personalizadas (P03).
  \item Anexar documentos para orientar a IA (P02).
  \item Funcionamento \emph{offline} (T08).
  \item Recomendações pedagógicas por IA e pré-visualização (V1-02 a V1-04:
    são compromissos, não funções existentes).
\end{itemize}
\end{limitacao}

\section{Decisões que o TR precisa tomar}

\begin{decisao}
\begin{enumerate}
  \item \textbf{Residência de dados}: admitir transferência internacional com
    as salvaguardas da ANPD ou exigir dados no Brasil (T06). A segunda opção
    muda infraestrutura, custo e o modelo de IA disponível.
  \item \textbf{Retenção e exclusão}: prazos por categoria de dado e o
    destino ao fim do contrato (T06).
  \item \textbf{Revogação de acesso}: prazo máximo entre desligar uma pessoa e
    ela perder acesso, inclusive em sessão aberta (T02).
  \item \textbf{Integração cadastral}: se basta a rede enviar as cargas ou se
    será exigida sincronização automática (T03).
  \item \textbf{Nível de serviço e suporte}: disponibilidade, horários,
    severidades e prazos (T09, E03).
  \item \textbf{Equipamento de referência}: navegadores, telas e banda mínimos
    a declarar no edital (T01).
\end{enumerate}
\end{decisao}

\section{Pendências por time}

\begin{center}
\footnotesize\renewcommand{\arraystretch}{1.35}
\begin{tabularx}{\linewidth}{@{}L{34mm}X@{}}
\toprule
\textbf{Time} & \textbf{O que precisa entregar para fechar o TR}\\
\midrule
\endhead
Produto / Pedagogia & Relatório da PoC (P10); capturas de ficha e guia (P01);
  par entrada/saída de personalização (P02); prompts e casos de avaliação
  das recomendações (V1-03, V1-04); responsáveis e datas da V1.\\
Editorial & Matriz ano $\times$ bimestre com livro e missão; componentes
  entregues; o que está no livro e o que está na plataforma (E01, E02).\\
Engenharia & Matriz de compatibilidade (T01); teste de carga do
  \cd{PU-EST-001} (T09); RPO e RTO; diagrama geral de fluxo de dados (T06);
  relatório atualizado de acessibilidade e avaliação manual (P09);
  configuração dos filtros do provedor de IA (T05); auditoria de
  autenticação e publicação (T07); V1-01, V1-02 e V1-06.\\
Jurídico / DPO & Retenção e exclusão; suboperadores (incluir o avatar da
  Ajuda); cláusula do provedor de IA sobre uso de dados (T06).\\
Implantação / Comercial & Plano de implantação, formação e suporte (E03,
  T09).\\
\bottomrule
\end{tabularx}
\end{center}

\section{Recomendações de forma para o TR}

\begin{itemize}
  \item Escrever cada requisito como resultado observável com critério de
    aceite, como nas caixas ``Proposta de redação'' deste documento, sem
    nomear tecnologia, nuvem ou marca.
  \item Exigir demonstração em homologação com dados sintéticos para todo
    requisito declarado como atendido.
  \item Tratar acessibilidade e proteção de dados como exigência legal, não
    como item de \emph{roadmap}.
  \item Manter fora do escopo mínimo as ideias que não são compromisso da V1
    (duplicação e compartilhamento, anexos, modo \emph{offline}, SAML) ou
    marcá-las como opcionais pontuáveis.
\end{itemize}
